\section{Pruebas y Validación}

\subsection{Estrategia de Pruebas}

Para garantizar la calidad y funcionamiento correcto del sistema, implementamos una estrategia de pruebas multi-nivel:

\begin{itemize}
    \item \textbf{Pruebas unitarias}: Tests de funciones y métodos individuales
    \item \textbf{Pruebas de integración}: Tests de interacción entre componentes
    \item \textbf{Pruebas de API}: Validación de endpoints con Swagger
    \item \textbf{Pruebas de carga}: Simulación de usuarios concurrentes
    \item \textbf{Pruebas de seguridad}: Validación de vulnerabilidades comunes
    \item \textbf{Pruebas manuales}: Testing exploratorio por usuarios reales
\end{itemize}

\subsection{Pruebas Backend (Django)}

\subsubsection{Testing Unitario}

Implementamos tests unitarios para los servicios críticos del sistema:

\begin{itemize}
    \item \textbf{Módulo Security}: Tests de autenticación 2FA, generación de tokens JWT, validación de permisos
    \item \textbf{Módulo General}: Tests de lógica de instructores (límites de aprendices, validación de contratos)
    \item \textbf{Módulo Assign}: Tests de flujo de asignación, cambios de estado, validaciones de fechas
\end{itemize}

\textbf{Herramientas}:
\begin{itemize}
    \item Django TestCase para tests de modelos
    \item APITestCase para tests de endpoints
    \item Mock para simular dependencias externas
\end{itemize}

\subsubsection{Testing de API con Swagger}

Utilizamos \textbf{drf-yasg} para generar documentación interactiva y testing manual:

\begin{itemize}
    \item Documentación automática de todos los endpoints
    \item Validación de esquemas de entrada y salida
    \item Testing manual directo desde el navegador
    \item Generación de código de cliente automático
\end{itemize}

URL de pruebas: \texttt{http://localhost:8000/swagger/}

\subsection{Pruebas Frontend (React + TypeScript)}

\subsubsection{Testing Unitario con Jest}

Configuramos Jest con React Testing Library para tests de componentes:

\begin{itemize}
    \item Tests de componentes de Login (6 archivos de test)
    \item Tests de Dashboard (2 archivos de test)
    \item Tests de ApplicationEvaluation
    \item Tests de hooks personalizados (2 archivos de test)
\end{itemize}

\textbf{Configuración}:
\begin{itemize}
    \item Jest 29.7 con jsdom environment
    \item React Testing Library para queries semánticas
    \item MSW (Mock Service Worker) para mockear APIs
    \item Coverage reports automáticos
\end{itemize}

\subsubsection{Validación de Tipos con TypeScript}

TypeScript 5.8 nos proporciona validación en tiempo de compilación:

\begin{itemize}
    \item Interfaces para todas las entidades (14 tipos)
    \item Tipado estricto de props de componentes
    \item Validación de respuestas de API
    \item Prevención de errores de null/undefined
\end{itemize}

\subsection{Pruebas de Integración}

\subsubsection{Flujos End-to-End Validados}

Validamos los flujos completos del sistema:

\begin{enumerate}
    \item \textbf{Flujo de Registro y Login}:
    \begin{itemize}
        \item Usuario ingresa credenciales
        \item Sistema envía código 2FA por correo
        \item Usuario ingresa código
        \item Sistema genera JWT y redirige según rol
    \end{itemize}
    
    \item \textbf{Flujo de Creación de Solicitud}:
    \begin{itemize}
        \item Aprendiz llena formulario de solicitud
        \item Sistema valida datos de empresa, jefe, talento humano
        \item Sistema guarda solicitud en estado SIN\_ASIGNAR
        \item Sistema envía notificación a coordinador
    \end{itemize}
    
    \item \textbf{Flujo de Asignación de Instructor}:
    \begin{itemize}
        \item Coordinador selecciona solicitud
        \item Sistema muestra instructores disponibles con límites
        \item Coordinador asigna instructor
        \item Sistema cambia estado a ASIGNADO
        \item Sistema genera 3 visitas automáticamente
        \item Sistema envía correo al instructor
    \end{itemize}
    
    \item \textbf{Flujo de Valoración}:
    \begin{itemize}
        \item Instructor revisa solicitud
        \item Instructor aprueba o rechaza con mensaje
        \item Sistema cambia estado a PRE-APROBADO o RECHAZADO
        \item Coordinador revisa y aprueba finalmente
        \item Sistema cambia estado a FINALIZADA
    \end{itemize}
\end{enumerate}

\subsection{Pruebas de Carga}

Utilizamos \textbf{Apache JMeter} para simular carga en el sistema:

\textbf{Escenarios de prueba}:

\begin{itemize}
    \item \textbf{Test 1: Usuarios concurrentes}
    \begin{itemize}
        \item 100 usuarios simultáneos navegando el sistema
        \item Duración: 10 minutos
        \item Resultado: Sistema estable, tiempos de respuesta < 500ms
    \end{itemize}
    
    \item \textbf{Test 2: Creación masiva de solicitudes}
    \begin{itemize}
        \item 500 solicitudes creadas en 5 minutos
        \item Resultado: BD manejó la carga sin problemas
    \end{itemize}
    
    \item \textbf{Test 3: Consultas concurrentes}
    \begin{itemize}
        \item 200 usuarios consultando listados simultáneamente
        \item Resultado: Redis caché mejoró rendimiento en 60\%
    \end{itemize}
    
    \item \textbf{Test 4: Estrés máximo}
    \begin{itemize}
        \item 5000 usuarios concurrentes (escenario extremo)
        \item Resultado: Sistema aguantó hasta 3000 usuarios, después degradación
    \end{itemize}
\end{itemize}

\subsection{Pruebas de Seguridad}

Validamos las vulnerabilidades comunes del \textbf{OWASP Top 10}:

\begin{enumerate}
    \item \textbf{Inyección SQL}: Django ORM previene automáticamente
    \item \textbf{Autenticación rota}: JWT + 2FA + expiración de tokens
    \item \textbf{Exposición de datos sensibles}: Passwords hasheados con bcrypt, HTTPS en producción
    \item \textbf{XXE (XML External Entities)}: No usamos XML
    \item \textbf{Control de acceso roto}: Sistema de permisos granular validado en cada endpoint
    \item \textbf{Configuración incorrecta}: Variables de entorno, \texttt{DEBUG=False} en producción
    \item \textbf{XSS (Cross-Site Scripting)}: Django escapa HTML automáticamente, React también
    \item \textbf{Deserialización insegura}: Validación con Serializers
    \item \textbf{Componentes vulnerables}: Actualizaciones regulares de dependencias
    \item \textbf{Logging insuficiente}: Logs con django.log y monitoreo de errores
\end{enumerate}

\subsection{Validación con Usuarios Reales}

Realizamos pruebas con usuarios reales del SENA:

\textbf{Participantes}:
\begin{itemize}
    \item 3 coordinadores de etapa productiva
    \item 5 instructores de seguimiento
    \item 10 aprendices en etapa productiva
\end{itemize}

\textbf{Resultados}:
\begin{itemize}
    \item \textbf{Facilidad de uso}: 8.5/10 promedio
    \item \textbf{Claridad de interfaz}: 9/10 promedio
    \item \textbf{Velocidad del sistema}: 8/10 promedio
    \item \textbf{Bugs encontrados}: 12 bugs menores corregidos
\end{itemize}

\subsection{Métricas de Calidad de Código}

Utilizamos herramientas automatizadas para medir calidad:

\textbf{Backend (Python)}:
\begin{itemize}
    \item \textbf{Pylint}: Score de 8.5/10
    \item \textbf{Complejidad ciclomática}: Promedio de 4.2 (bueno)
    \item \textbf{Cobertura de tests}: 45\% (objetivo: 70\%)
\end{itemize}

\textbf{Frontend (TypeScript)}:
\begin{itemize}
    \item \textbf{ESLint}: 23 warnings menores, 0 errors
    \item \textbf{TypeScript strict mode}: Habilitado
    \item \textbf{Cobertura de tests}: 35\% (en crecimiento)
\end{itemize}

\subsection{Lecciones de las Pruebas}

Las pruebas nos revelaron áreas importantes de mejora:

\begin{enumerate}
    \item \textbf{Optimización de consultas}: Detectamos N+1 queries que se resolvieron con \texttt{select\_related}
    \item \textbf{Validaciones del lado del cliente}: Mejorar feedback visual antes de enviar al servidor
    \item \textbf{Manejo de errores}: Agregar más mensajes descriptivos para el usuario
    \item \textbf{Performance en móviles}: Optimizar para dispositivos con menos recursos
\end{enumerate}

